% Standalone problem document, generated from the adjacent README.md.
% Compile with XeLaTeX. Mathematical content is maintained in Markdown.
\documentclass[11pt,a4paper]{article}
\usepackage[fontsize=10.5pt]{scrextend}
\usepackage[margin=22mm,headheight=15pt,headsep=9mm,footskip=11mm]{geometry}
\usepackage{fontspec}
\setmainfont{texgyrepagella}[Extension=.otf,UprightFont=*-regular,BoldFont=*-bold,ItalicFont=*-italic,BoldItalicFont=*-bolditalic]
\setsansfont{texgyreheros}[Extension=.otf,UprightFont=*-regular,BoldFont=*-bold,ItalicFont=*-italic,BoldItalicFont=*-bolditalic]
\setmonofont{texgyrecursor}[Extension=.otf,UprightFont=*-regular,BoldFont=*-bold,ItalicFont=*-italic,BoldItalicFont=*-bolditalic,Scale=MatchLowercase]
\usepackage{amsmath,amssymb}
\usepackage{unicode-math}
\setmathfont{texgyrepagella-math.otf}
\usepackage{xcolor,microtype,enumitem,needspace,fancyhdr}
\usepackage{bookmark}
\definecolor{ink}{HTML}{17344B}
\definecolor{accent}{HTML}{197D80}
\definecolor{muted}{HTML}{596773}
\hypersetup{colorlinks=true,linkcolor=accent,urlcolor=accent,pdftitle={MI-11 - Lieb's
permanental dominance
conjecture},pdfauthor={Open Problems in Numerical Linear Algebra}}
\urlstyle{same}
\setlength{\parindent}{0pt}
\setlength{\parskip}{5pt plus 1pt minus 1pt}
\linespread{1.04}
\setlength{\emergencystretch}{3em}
\widowpenalty=10000
\clubpenalty=10000
\displaywidowpenalty=10000
\allowdisplaybreaks[1]
\setlist{nosep,leftmargin=1.5em}
\providecommand{\tightlist}{\setlength{\itemsep}{0pt}\setlength{\parskip}{0pt}}
\providecommand{\pandocbounded}[1]{#1}
\pagestyle{fancy}
\fancyhf{}
\fancyhead[L]{\sffamily\footnotesize\color{muted}OPEN PROBLEMS IN NUMERICAL LINEAR ALGEBRA}
\fancyhead[R]{\sffamily\bfseries\footnotesize\color{accent}MI-11}
\fancyfoot[L]{\parbox[b]{0.9\textwidth}{\sffamily\fontsize{8}{10}\selectfont\color{muted}Editorial ratings; a bounded search cannot certify that a problem remains unsolved.\\Literature check: 2026-09-08}}
\fancyfoot[R]{\sffamily\footnotesize\color{muted}\thepage}
\renewcommand{\headrulewidth}{0.3pt}
\renewcommand{\headrule}{\hbox to\headwidth{\color{accent}\leaders\hrule height \headrulewidth\hfill}}
\makeatletter
\renewcommand{\section}{\Needspace{5\baselineskip}\@startsection{section}{1}{0pt}{12pt plus 2pt minus 2pt}{4pt}{\sffamily\bfseries\large\color{ink}}}
\renewcommand{\subsection}{\Needspace{4\baselineskip}\@startsection{subsection}{2}{0pt}{9pt plus 1pt minus 1pt}{3pt}{\sffamily\bfseries\normalsize\color{ink}}}
\makeatother
\setcounter{secnumdepth}{0}
\begin{document}
{\sffamily\small\color{muted}Matrix inequalities and norms\par}
\vspace{3pt}
{\raggedright\hyphenpenalty=10000\exhyphenpenalty=10000\sffamily\bfseries\fontsize{21}{24}\selectfont\color{ink}Lieb's
permanental dominance conjecture\par}
\vspace{7pt}
{\sffamily\small\textbf{Difficulty:} extreme\quad\textbf{Importance:} interesting
to the community\par}
\vspace{4pt}
{\color{accent}\hrule height 0.8pt}
\vspace{6pt}
\textbf{Status:} explicit conjecture; no general resolution located

Let \(n\ge1\), let \(A=(a_{ij})\in\mathbb C^{n\times n}\) be Hermitian
positive semidefinite, let \(G\) be any subgroup of the symmetric group
\(S_n\), and let \(\chi\) be the character of any nonzero
finite-dimensional complex representation of \(G\). Thus \(\chi(e)>0\),
where \(e\) is the identity. Define \[
f_\chi(A)=\frac{1}{\chi(e)}
\sum_{\sigma\in G}\chi(\sigma)\prod_{i=1}^n a_{i,\sigma(i)},
\qquad
\operatorname{per}A=\sum_{\sigma\in S_n}\prod_{i=1}^n a_{i,\sigma(i)}.
\] Does the inequality \[
f_\chi(A)\le\operatorname{per}A
\] hold for every such \(n,A,G,\chi\)? The character expression is real
for Hermitian \(A\), so the comparison is an ordinary real inequality.

\subsection{Relevance}\label{relevance}

Generalized matrix functions include the determinant and immanants. The
conjecture seeks a common upper bound across these structured matrix
polynomials on the positive-semidefinite cone. It is part of matrix
computation and matrix inequality theory, with links to symmetry
reductions of tensor powers.

\subsection{References}\label{references}

\begin{itemize}
\tightlist
\item
  E. H. Lieb, \emph{Proofs of some Conjectures on Permanents}, Journal
  of Mathematics and Mechanics 16 (1966), 127--134, Conjecture
  \(\alpha\) on p.127
  (\href{https://iumj.s3-us-west-2.amazonaws.com/abstracts/16008_abs.pdf}{publisher's
  first page}).
\item
  I. M. Wanless, \emph{Lieb's permanental dominance conjecture} (2022),
  Conjecture 1, §2 implication diagram, and §3 partial results
  (\href{https://arxiv.org/pdf/2202.01867}{primary manuscript}).
\item
  A. Rico, D. Grinko, R. Krebs, and L. H. Zaw, \emph{Entanglement
  Structure and Matrix Inequalities from Isotypic Measurements},
  Physical Review Letters 137 (2026), 100203; abstract describes
  immanant inequalities for orders three and four
  (\href{https://doi.org/10.1103/nvk2-h8d5}{journal}).
\end{itemize}

\subsection{Status check --- 2026-09-08}\label{status-check-2026-09-08}

Searches for the exact conjecture name and ``proof'',
``counterexample'', and 2025--2026 found no general resolution. The
September 2026 PRL abstract describes fixed-order progress. The August
2026 preprint \emph{GPT, the Counterexample Machine}, §3.1, refutes a
stronger conjecture for arbitrary real-stable polynomials; its example
is not presented as a PSD-matrix counterexample to Lieb's conjecture
(\href{https://arxiv.org/html/2608.29595v1}{primary manuscript}).
Likewise, the older disproof of permanent-on-top does not resolve the
displayed statement.
\end{document}
